% !TeX root = ../thuthesis-example.tex
%
% 第三章：路由算法与下界分析
%

\chapter{路由算法与下界分析}
\label{chap:nonisotonic}

本章在第~\ref{chap:phymodel}章路径级物理模型的基础上，依次讨论量子路由面对的主要困难、单标签 Dijkstra 基线、路径级保真度度量的非保序性、单标签 Dijkstra 相对于实际最优路径的下界、切片 Dijkstra 的提出，以及切片算法对应的下界保证。本文把“实际最优路径”定义为在同一网络拓扑、同一链路初始保真度、同一建链时间和同一相干时间参数下，使 $F_{\mathrm{path}}(p)$ 最大的简单路径。

\section{量子路由的困难}
\label{sec:quantum-routing-difficulty}

设网络图为 $G=(V,E)$，源节点为 $s$，目的节点为 $t$，从 $s$ 到 $t$ 的简单路径集合记为 $\mathcal{P}_{s,t}$。对任意路径
\[
p=(v_0,v_1,\ldots,v_h),
\qquad v_0=s,\quad v_h=t,
\]
路径级保真度由第~\ref{chap:phymodel}章的等待退相干和 Bell 对角态交换递推给出，记为 $F_{\mathrm{path}}(p)$。实际最优路径为
\begin{equation}
  p^\star
  \in
  \arg\max_{p\in\mathcal{P}_{s,t}} F_{\mathrm{path}}(p),
  \qquad
  F^\star = F_{\mathrm{path}}(p^\star).
  \label{eq:actual-optimal-path}
\end{equation}

与经典加性最短路不同，本文的路径质量函数具有三个直接后果。第一，新增一段链路会改变整条路径的 $T_{\max}$，从而改变已有链路的等待时间；第二，等待后的各段 Bell 对角态再经过递归纠缠交换，端到端保真度不能由单条边权相加得到；第三，中间节点处较优的前缀，在接入同一后缀后未必仍然较优。因此，路由算法的核心困难在于搜索过程中如何保留未来可能反超的前缀，并避免因局部比较过早剪枝。

\section{单标签 Dijkstra 基线}
\label{sec:dijkstra-baseline}

最直接的基线是单标签 Dijkstra。该算法对每个节点 $v$ 只保留一条从 $s$ 到 $v$ 的候选前缀，记为 $\ell(v)$。优先队列按照候选前缀的 $F_{\mathrm{path}}$ 从高到低弹出；若两个候选的保真度相同，则优先选择 $T_{\max}$ 更小、路径更短的候选。每次从队列中取出一个前缀后，算法将其末端节点的出边逐一接入，并用新路径的完整物理参数重新计算 $F_{\mathrm{path}}$ 和 $T_{\max}$。

单标签 Dijkstra 的计算代价较低。若忽略路径级重算的内部代价，它与普通最好优先搜索类似；若一条长度为 $h$ 的前缀每次扩展都需要重新计算等待退相干和 Bell 对角态交换，则一次扩展的物理度量计算代价为 $O(h)$。该算法的关键限制在于：每个中间节点只有一个保留前缀。只要被删除的前缀在接入后缀后可能反超，单标签剪枝就可能破坏全局最优性。

\section{非保序性的论证}
\label{sec:non-isotonicity}

\subsection{保序性的定义}

若一个路径度量 $F$ 对任意两条同终点前缀 $a,b$ 和任意可拼接后缀 $c$ 满足
\begin{equation}
  F(a) \ge F(b)
  \implies
  F(a\oplus c) \ge F(b\oplus c),
  \label{eq:isotonicity-def}
\end{equation}
则称该度量在路径拼接下是保序的。经典 Dijkstra 的局部剪枝依赖这一结构：一旦某个前缀在中间节点处已经最优，拼接相同后缀后它仍然不劣。

本文的 $F_{\mathrm{path}}$ 不满足式~\eqref{eq:isotonicity-def}。原因在于，路径拼接会重新决定 $T_{\max}$，并由此重新分配所有链路的等待时间。一个前缀在中间节点处保真度较高，可能同时携带较大的时间瓶颈；当它接入一个快速后缀时，后缀需要等待该时间瓶颈，从而发生额外退相干。另一个局部保真度稍低但时间瓶颈更小的前缀，则可能在同一后缀下保持更高的最终保真度。

\subsection{构造性反例}

先考虑一条单段 Werner 路径。若其初始保真度为 $F$，左右端相干时间均为 $\tau$，并且交换前等待 $\Delta$，则由式~\eqref{eq:single-link-fidelity} 可得等待后的 $\Phi^+$ 权重为
\begin{equation}
  g(F,\Delta;\tau)
  =
  \frac{2F + 1 + e^{-2\Delta/\tau}(4F - 1)}{6}.
  \label{eq:waited-werner-fidelity}
\end{equation}
在 $F>1/4$ 的物理区间内，$g$ 随 $F$ 单调递增，随 $\Delta$ 单调递减。

两段 Werner 态经过理想纠缠交换后的 $\Phi^+$ 保真度为
\begin{equation}
  S(F_1,F_2)
  =
  F_1F_2 + \frac{(1-F_1)(1-F_2)}{3}.
  \label{eq:werner-swap}
\end{equation}
当 $F_1,F_2>1/4$ 时，$S$ 关于每个自变量单调递增，因为
\[
  \frac{\partial S}{\partial F_1}
  =
  \frac{4F_2 - 1}{3},
  \qquad
  \frac{\partial S}{\partial F_2}
  =
  \frac{4F_1 - 1}{3}.
\]

取两个到同一中间节点的前缀 $a,b$ 和一个共享后缀 $x$。令
\[
F_{\mathrm{path}}(a)>F_{\mathrm{path}}(b),
\qquad
T_{\max}(a)>T_{\max}(b).
\]
若后缀 $x$ 的建链时间较小，则在 $a\oplus x$ 中，$x$ 需要等待的时间可能大于其在 $b\oplus x$ 中的等待时间。只要
\[
S\!\bigl(F_{\mathrm{path}}(a),g(F_x,\Delta_x^{(a)};\tau)\bigr)
<
S\!\bigl(F_{\mathrm{path}}(b),g(F_x,\Delta_x^{(b)};\tau)\bigr),
\qquad
\Delta_x^{(a)}>\Delta_x^{(b)},
\]
就得到
\[
F_{\mathrm{path}}(a\oplus x)
<
F_{\mathrm{path}}(b\oplus x).
\]

例如取
\[
F_a=0.55,\quad T_a=5\ \mathrm{ms},
\qquad
F_b=0.50,\quad T_b=0\ \mathrm{ms},
\]
\[
F_x=0.90,\quad T_x=0\ \mathrm{ms},
\qquad
\tau=20\ \mathrm{ms}.
\]
不接后缀时有 $F_{\mathrm{path}}(a)=0.55>0.50=F_{\mathrm{path}}(b)$。接入同一后缀后，
\[
g(0.90,5;20)\approx 0.7295,
\qquad
g(0.90,0;20)=0.90.
\]
代入式~\eqref{eq:werner-swap} 得
\[
F_{\mathrm{path}}(a\oplus x)\approx 0.4418,
\qquad
F_{\mathrm{path}}(b\oplus x)\approx 0.4667.
\]
因此，路径级保真度在实际等待模型下发生排序反转，单标签 Dijkstra 的局部剪枝前提不成立。

\section{单标签 Dijkstra 的下界证明}
\label{sec:dijkstra-lower-bound}

非保序性说明单标签 Dijkstra 不能保证找到 $p^\star$。本节进一步分析它相对于实际最优路径 $p^\star$ 的下界。结论分为两部分：其一，在给定实际参数范围后，任何返回的可行路径都有一个可计算的保真度下界；其二，若不加入额外保序假设，单标签 Dijkstra 相对于 $F^\star$ 的最坏情形只能保证这一平凡下界。

令可用链路集合为 $E_{\mathrm{a}}$。定义全图建链时间跨度
\begin{equation}
  \Delta_G
  =
  \max_{e\in E_{\mathrm{a}}} T(e)
  -
  \min_{e\in E_{\mathrm{a}}} T(e),
  \label{eq:global-time-span}
\end{equation}
并令 $H=|V|-1$ 为简单路径的最大跳数。对一条链路 $e=(u,v)$，记其初始 Werner 保真度为 $F_e^{(0)}$，两端相干时间为 $\tau_u,\tau_v$。在等待时间 $w$ 下，其等待后的 $\Phi^+$ 权重为
\begin{equation}
  a_e(w)
  =
  \frac{
    2F_e^{(0)}+1+
    \exp\!\left(-\frac{w}{\tau_u}-\frac{w}{\tau_v}\right)
    (4F_e^{(0)}-1)
  }{6}.
  \label{eq:edge-wait-lower-component}
\end{equation}
定义由实际参数诱导的单链路保真度下界
\begin{equation}
  \alpha_G
  =
  \min_{e\in E_{\mathrm{a}}}\;
  \min_{0\le w\le \Delta_G}
  a_e(w).
  \label{eq:alpha-g}
\end{equation}
若所有可用链路满足 $F_e^{(0)}\ge 1/4$，则 $a_e(w)$ 关于 $w$ 单调不增，此时内层最小值在 $w=\Delta_G$ 处取得。

\paragraph{命题 1}
对任意简单路径 $p$，若其跳数为 $h(p)$，则
\begin{equation}
  F_{\mathrm{path}}(p)
  \ge
  \alpha_G^{h(p)}
  \ge
  \alpha_G^{H}.
  \label{eq:any-path-lower-bound}
\end{equation}

\paragraph{证明}
路径 $p$ 上任意链路的等待时间均不超过 $\Delta_G$，故等待后该链路 Bell 对角态的第一个分量不小于 $\alpha_G$。设两个 Bell 对角态的权重向量分别为 $w=(a,b,c,d)$ 与 $w'=(a',b',c',d')$。由式~\eqref{eq:bell-swap-operator}，交换后第一个分量为
\[
aa' + bb' + cc' + dd'
\ge aa'.
\]
因此，每进行一次交换，输出态的第一个分量至少为两个输入第一个分量的乘积。沿 $h(p)$ 条链路递推，即得
\[
F_{\mathrm{path}}(p)\ge \prod_{e\in p}\alpha_G=\alpha_G^{h(p)}.
\]
又因 $h(p)\le H$ 且 $0\le\alpha_G\le 1$，故式~\eqref{eq:any-path-lower-bound} 成立。 \(\square\)

设单标签 Dijkstra 返回路径 $p_D$，跳数为 $h_D$，其保真度为 $F_D=F_{\mathrm{path}}(p_D)$。由于 $F^\star\le 1$，由命题 1 立刻得到相对于实际最优路径的下界
\begin{equation}
  \frac{F_D}{F^\star}
  \ge
  F_D
  \ge
  \alpha_G^{h_D}
  \ge
  \alpha_G^H.
  \label{eq:dijkstra-lower-bound}
\end{equation}
等价地，加性损失满足
\begin{equation}
  F^\star - F_D
  \le
  1-\alpha_G^{h_D}
  \le
  1-\alpha_G^H.
  \label{eq:dijkstra-additive-bound}
\end{equation}

式~\eqref{eq:dijkstra-lower-bound} 是一个由实际网络参数计算得到的严格下界，但它通常较弱。其原因在于单标签 Dijkstra 在非保序度量下确实可能删除实际最优路径的必要前缀。对任意给定的中间节点，可以构造两个前缀 $a,b$ 和共享后缀 $x$，使 $F_{\mathrm{path}}(a)>F_{\mathrm{path}}(b)$，但 $F_{\mathrm{path}}(a\oplus x)<F_{\mathrm{path}}(b\oplus x)$。单标签 Dijkstra 会保留 $a$ 并删除 $b$，从而无法生成包含 $b\oplus x$ 的最优路径。因此，在不增加额外结构假设的情况下，不能证明单标签 Dijkstra 具有比式~\eqref{eq:dijkstra-lower-bound} 更强的普适近似保证。

\section{切片 Dijkstra 算法}
\label{sec:sliced-dijkstra}

上述失败机制表明，单个中间节点处只保留一条前缀过于激进。切片 Dijkstra 的基本思想是沿 $T_{\max}$ 轴保留多个代表性前缀，使时间瓶颈不同的候选能够共存。给定桶宽 $\Delta_B>0$，定义路径 $p$ 的时间切片编号为
\begin{equation}
  \beta(p)
  =
  \left\lfloor
    \frac{T_{\max}(p)}{\Delta_B}
  \right\rfloor .
  \label{eq:time-slice-id}
\end{equation}
算法在每个节点维护至多 $K$ 条候选前缀。不同时间切片中的候选可以同时保留；同一节点、同一时间切片内，优先保留 $F_{\mathrm{path}}$ 更高的候选，若保真度相同，则选择 $T_{\max}$ 更小、跳数更少的候选。

\begin{algorithm}[htb]
  \caption{切片 Dijkstra 的多标签搜索流程}
  \label{alg:sliced-dijkstra}
  \begin{algorithmic}[1]
    \REQUIRE 源节点 $s$，目的节点 $t$，拓扑 $G$，每节点候选上限 $K$，桶宽 $\Delta_B$
    \ENSURE 从 $s$ 到 $t$ 的候选路径
    \STATE $\mathcal{L}(v) \gets \emptyset,\ \forall v\in V$
    \STATE 构造以 $s$ 为末端的初始前缀 $\ell_0$，并置 $\mathcal{L}(s)\gets\{\ell_0\}$
    \STATE 将 $\ell_0$ 插入按 $F_{\mathrm{path}}$ 排序的优先队列 $Q$
    \STATE $p_{\mathrm{best}}\gets\emptyset$
    \WHILE{$Q$ 非空}
      \STATE 取出 $Q$ 中保真度最高的候选前缀 $\ell$
      \STATE $u\gets \ell$ 的末端节点
      \IF{$\ell\notin \mathcal{L}(u)$}
        \STATE \textbf{continue}
      \ENDIF
      \IF{$u=t$}
        \STATE 若 $\ell$ 优于 $p_{\mathrm{best}}$，则 $p_{\mathrm{best}}\gets \ell$
        \STATE \textbf{continue}
      \ENDIF
      \FOR{每条可用出边 $(u,v)$}
        \STATE 令 $\ell'$ 为 $\ell$ 接入 $(u,v)$ 后得到的新前缀
        \STATE 根据时间切片和候选上限更新 $\mathcal{L}(v)$
        \IF{$\ell'$ 被 $\mathcal{L}(v)$ 保留}
          \STATE 将 $\ell'$ 插入 $Q$
        \ENDIF
      \ENDFOR
    \ENDWHILE
    \RETURN $p_{\mathrm{best}}$
  \end{algorithmic}
\end{algorithm}

设每个节点最多保留 $K$ 条候选前缀，则活跃前缀数为 $O(K|V|)$，边扩展次数为 $O(K|E|)$。每次扩展若对长度为 $h$ 的前缀重新计算 $T_{\max}$、等待退相干和 Bell 对角态交换，则物理度量计算代价为 $O(h)$。因此，切片算法用 $K$ 倍级别的候选保留开销，换取对不同时间瓶颈前缀的覆盖。

\section{切片算法的下界证明}
\label{sec:sliced-lower-bound}

切片 Dijkstra 的下界可以写得比单标签 Dijkstra 更细。关键不只是它返回了一条可行路径，还在于它实际保留下了一部分完整候选路径。令 $\mathcal{R}_{K,\Delta_B}\subseteq \mathcal{P}_{s,t}$ 表示切片算法在给定 $K$ 与 $\Delta_B$ 下能够完整生成并保留到终点的路径集合。算法返回路径 $p_S$ 满足
\begin{equation}
  F_S
  =
  F_{\mathrm{path}}(p_S)
  =
  \max_{p\in\mathcal{R}_{K,\Delta_B}}
  F_{\mathrm{path}}(p).
  \label{eq:sliced-retained-best}
\end{equation}
因此，相对于实际最优路径有
\begin{equation}
  \frac{F_S}{F^\star}
  \ge
  \frac{
    \max_{p\in\mathcal{R}_{K,\Delta_B}}F_{\mathrm{path}}(p)
  }{
    F^\star
  }.
  \label{eq:sliced-retained-lower-bound}
\end{equation}
式~\eqref{eq:sliced-retained-lower-bound} 是切片算法比单标签基线更有意义的实例级下界：它直接由算法保留下来的完整路径集合决定。

进一步，由命题 1 可知，只要切片算法返回一条可行简单路径 $p_S$，就有
\begin{equation}
  \frac{F_S}{F^\star}
  \ge
  \alpha_G^{h(p_S)}
  \ge
  \alpha_G^H.
  \label{eq:sliced-unconditional-lower-bound}
\end{equation}
这给出与单标签 Dijkstra 相同形式的无条件参数下界。

切片算法能够取得更强保证的条件可以形式化为前缀覆盖。设 $p^\star=(v_0,\ldots,v_h)$，其第 $i$ 个前缀记为
\[
p^\star_i=(v_0,\ldots,v_i).
\]
若对所有 $i=0,\ldots,h$，前缀 $p^\star_i$ 在算法处理节点 $v_i$ 时均被保留，则 $p^\star\in\mathcal{R}_{K,\Delta_B}$，从而由式~\eqref{eq:sliced-retained-best} 得
\begin{equation}
  F_S = F^\star.
  \label{eq:sliced-exact-condition}
\end{equation}
因此，切片 Dijkstra 的精确性不来自时间切片本身，而来自它是否保留了实际最优路径的所有必要前缀。

需要强调的是，有限 $K$ 和有限桶宽并不能自动保证式~\eqref{eq:sliced-exact-condition} 的前缀覆盖条件。若实际最优路径的某个前缀与另一个局部保真度更高的前缀落入同一时间桶，且该桶只保留一个代表，则最优前缀仍可能被删除。此时切片算法不一定优于单标签 Dijkstra，其无条件最坏情形下界仍为式~\eqref{eq:sliced-unconditional-lower-bound}。换言之，切片算法提高的是保留关键前缀的机会；严格最优性需要由前缀覆盖条件或全路径枚举来保证。
